\section{Лекция 4. Формула бинома Ньютона и биномиальные коэффициенты}
\noeqref{eq:newton_binom, eq:binom_sym, eq:binom_sum_all, eq:binom_sum_alt}
\noeqref{eq:floor_ceil_def, eq:binom_sum_even_odd, eq:binom_sum_k}
\noeqref{eq:binom_sum_int, eq:pascal_id, eq:vandermonde, eq:hockey_stick}
\noeqref{eq:fib_binom, eq:prime_binom_crit}

\subsection{Полиномы и формула бинома Ньютона}

\begin{definition}
    Введём базовые алгебраические понятия для выражений от одной или 
    нескольких переменных:
    \begin{itemize}
        \item \textdef{Моном} (\textdef{одночлен}) --- произведение 
        числового коэффициента и переменных в целых неотрицательных 
        степенях (например: $x$, $1$, $x y^2$).
        \item \textdef{Бином} (\textdef{двучлен}) --- алгебраическая сумма 
        двух мономов (например: $x + 1$, $x + y$).
        \item \textdef{Полином} (\textdef{многочлен}) --- конечная 
        алгебраическая сумма мономов.
    \end{itemize}
\end{definition}

\begin{theorem}[\textbf{Формула бинома Ньютона}]
    Для любых элементов $x$ и $y$ (из коммутативного кольца) и любого 
    целого неотрицательного показателя $n \in \N \cup \{0\}$ справедливо 
    разложение:
    \begin{equation} \label{eq:newton_binom}
        (x + y)^n = \sum_{k=0}^{n} C_n^k x^{n-k} y^k = 
        \sum_{k=0}^{n} \binom{n}{k} x^{n-k} y^k \tag{1}
    \end{equation}
    где числа $C_n^k = \binom{n}{k} = \frac{n!}{k!(n - k)!}$ называются 
    \textdef{биномиальными коэффициентами}.
\end{theorem}

\begin{proof}[Доказательство (Способ 1: комбинаторный вывод)]~\\
    \begin{quote}
        Представим $n$-ю степень двучлена $(x + y)^n$ в виде произведения 
        $n$ одинаковых скобок:
        \[
            (x + y)^n = \underbrace{(x + y)(x + y) \dots (x + y)}_{n \text{ раз}}
        \]
        По правилу раскрытия скобок каждое элементарное слагаемое в 
        результате перемножения образуется путём выбора ровно одного 
        множителя ($x$ или $y$) из каждой из $n$ скобок. В силу 
        коммутативности умножения произведение, в котором из $k$ скобок 
        была выбрана переменная $y$, а из остальных $n - k$ скобок --- 
        переменная $x$, приводится к виду монома $x^{n-k} y^k$.

        Следовательно, после приведения подобных слагаемых коэффициент при 
        одночлене $x^{n-k} y^k$ в точности равен количеству способов 
        выбрать $k$ скобок (из которых берётся множитель $y$) среди $n$ 
        имеющихся скобок без учёта порядка и без повторений. По определению 
        сочетаний это число равно $C_n^k$, что доказывает формулу 
        \eqref{eq:newton_binom}.
    \end{quote}
\end{proof}

\begin{proof}[Доказательство (Способ 2: метод математической индукции)]~\\
    \begin{quote}
        Проведём индукцию по показателю степени $n$.

        \textbf{1. База индукции ($n = 0$ и $n = 1$):} 
        При $n = 0$ левая часть равна $(x + y)^0 = 1$, а правая часть 
        состоит из одного слагаемого $C_0^0 x^0 y^0 = 1$. При $n = 1$ имеем 
        $(x + y)^1 = C_1^0 x^1 y^0 + C_1^1 x^0 y^1 = x + y$. База верна.

        \textbf{2. Индукционный переход ($n \implies n + 1$):} 
        Предположим, что формула \eqref{eq:newton_binom} верна для степени 
        $n$. Умножим обе части равенства \eqref{eq:newton_binom} на 
        $(x + y)$:
        \begin{align*}
            (x + y)^{n+1} &= (x + y) \sum_{k=0}^{n} C_n^k x^{n-k} y^k = 
            \sum_{k=0}^{n} C_n^k x^{n+1-k} y^k + 
            \sum_{k=0}^{n} C_n^k x^{n-k} y^{k+1}.
        \end{align*}
        Во второй сумме выполним сдвиг индекса суммирования $j = k + 1$ 
        (тогда $k = j - 1$, а $j$ пробегает значения от $1$ до $n + 1$), и 
        отделим крайние слагаемые при $k = 0$ и $j = n + 1$:
        \[
            (x + y)^{n+1} = C_n^0 x^{n+1} + 
            \sum_{k=1}^{n} \bigl(C_n^k + C_n^{k-1}\bigr) x^{n+1-k} y^k + 
            C_n^n y^{n+1}.
        \]
        Используя тождество Паскаля $C_n^k + C_n^{k-1} = C_{n+1}^k$ и 
        равенства $C_n^0 = C_{n+1}^0 = 1$, $C_n^n = C_{n+1}^{n+1} = 1$, 
        объединяем все слагаемые в единую сумму:
        \[
            (x + y)^{n+1} = \sum_{k=0}^{n+1} C_{n+1}^k x^{n+1-k} y^k.
        \]
        Индукционный переход доказан.
    \end{quote}
\end{proof}

\subsection{Свойства биномиальных коэффициентов}

\subsubsection{Свойство 1. Симметричность}

\begin{proposition}[\textbf{Симметричность биномиальных коэффициентов}]
    Для любых целых чисел $0 \le k \le n$ справедливо равенство:
    \begin{equation} \label{eq:binom_sym}
        C_n^k = C_n^{n-k} \tag{2}
    \end{equation}
\end{proposition}

\begin{proof}[Доказательство (алгебраическое и комбинаторное)]~\\
    \begin{quote}
        \textbf{1. Алгебраический вывод:} 
        Подставим индекс $n - k$ непосредственно в факториальную формулу 
        числа сочетаний:
        \[
            C_n^{n-k} = \frac{n!}{(n - k)!\bigl(n - (n - k)\bigr)!} = 
            \frac{n!}{(n - k)! k!} = \frac{n!}{k!(n - k)!} = C_n^k.
        \]

        \textbf{2. Комбинаторный вывод:} 
        Каждому выбору $k$-элементного подмножества $A$ из $n$-элементного 
        множества $X$ взаимно однозначно соответствует выбор его дополнения 
        $X \setminus A$, содержащего ровно $n - k$ элементов (выбрать $k$ 
        элементов, которые мы «берём», --- то же самое, что выбрать $n - k$ 
        элементов, которые мы «оставляем»). По принципу биективности 
        количества таких выборок равны: $C_n^k = C_n^{n-k}$.
    \end{quote}
\end{proof}

\subsubsection{Свойство 2. Сумма всех биномиальных коэффициентов}

\begin{proposition}
    Сумма всех биномиальных коэффициентов $n$-й строки равна $2^n$:
    \begin{equation} \label{eq:binom_sum_all}
        \sum_{k=0}^{n} C_n^k = 2^n \tag{3}
    \end{equation}
\end{proposition}

\begin{proof}[Доказательство]~\\
    \begin{quote}
        \textbf{1. По формуле бинома Ньютона:} 
        Положим в тождестве \eqref{eq:newton_binom} $x = 1$ и $y = 1$:
        \[
            (1 + 1)^n = \sum_{k=0}^{n} C_n^k \cdot 1^{n-k} \cdot 1^k \implies 
            2^n = \sum_{k=0}^{n} C_n^k.
        \]

        \textbf{2. Комбинаторное доказательство:} 
        Мощность булеана $\abs{\powerset(X)}$ множества из $n$ элементов 
        равна $2^n$. С другой стороны, все подмножества множества $X$ можно 
        разбить на $n + 1$ попарно непересекающихся классов по числу 
        элементов $k \in \{0, 1, \dots, n\}$. Число подмножеств мощности 
        $k$ равно $C_n^k$, откуда по принципу сложения $\sum_{k=0}^n C_n^k = 2^n$.
    \end{quote}
\end{proof}

\subsubsection{Свойство 3. Знакочередующаяся сумма}

\begin{proposition}
    Знакочередующаяся сумма биномиальных коэффициентов при любом натуральном 
    $n \ge 1$ равна нулю:
    \begin{equation} \label{eq:binom_sum_alt}
        \sum_{k=0}^{n} (-1)^k C_n^k = 0 \tag{4}
    \end{equation}
\end{proposition}

\begin{proof}[Доказательство]~\\
    \begin{quote}
        \textbf{1. По формуле бинома Ньютона:} 
        Подставим в формулу \eqref{eq:newton_binom} значения $x = 1$ и 
        $y = -1$ (при $n \ge 1$):
        \[
            \bigl(1 + (-1)\bigr)^n = \sum_{k=0}^{n} C_n^k \cdot 1^{n-k} \cdot (-1)^k 
            \implies 0 = \sum_{k=0}^{n} (-1)^k C_n^k.
        \]

        \textbf{2. Комбинаторное доказательство:} 
        Зафиксируем элемент $x_1 \in X$. Любое подмножество $A \subseteq X$, 
        не содержащее $x_1$, образует пару с подмножеством 
        $A \cup \{x_1\}$, содержащим $x_1$. В каждой такой паре мощности 
        множеств отличаются на $1$, поэтому одно из них имеет чётную 
        мощность, а другое --- нечётную. Это устанавливает биекцию между 
        семейством подмножеств чётной мощности и семейством подмножеств 
        нечётной мощности, откуда их количества равны, а разность равна $0$.
    \end{quote}
\end{proof}

\subsubsection{Свойство 4. Суммы по чётным и нечётным индексам}

\begin{definition}
    Для любого вещественного числа $x \in \R$ определяются функции 
    округления до ближайшего целого снизу и сверху:
    \begin{itemize}
        \item \textdef{Целая часть «снизу» («пол», floor):} 
        наибольшее целое число, не превосходящее $x$:
        \[
            \lfloor x \rfloor = \max \{y \in \Z \mid y \le x\}
        \]
        \item \textdef{Целая часть «сверху» («потолок», ceiling):} 
        наименьшее целое число, не меньшее $x$:
        \[
            \lceil x \rceil = \min \{y \in \Z \mid y \ge x\}
        \]
    \end{itemize}
\end{definition}

\begin{proposition}
    При любом $n \ge 1$ сумма биномиальных коэффициентов, стоящих на чётных 
    местах, равна сумме коэффициентов на нечётных местах и составляет 
    $2^{n-1}$ (при соглашении $C_n^m = 0$ для $m > n$):
    \begin{equation} \label{eq:binom_sum_even_odd}
        \sum_{k=0}^{\lfloor n/2 \rfloor} C_n^{2k} = 
        \sum_{k=0}^{\lfloor n/2 \rfloor} C_n^{2k+1} = 2^{n-1} \tag{6}
    \end{equation}
\end{proposition}

\begin{proof}[Доказательство]~\\
    \begin{quote}
        Воспользуемся свойствами 2 и 3. В тождестве \eqref{eq:binom_sum_alt} 
        слагаемые с чётными индексами входят со знаком «$+$», а с 
        нечётными --- со знаком «$-$». Сложим почленно равенства 
        \eqref{eq:binom_sum_all} и \eqref{eq:binom_sum_alt}:
        \[
            \sum_{k=0}^{n} C_n^k + \sum_{k=0}^{n} (-1)^k C_n^k = 
            2 \sum_{\substack{k=0 \\ k \text{ чётные}}}^{n} C_n^k = 2^n + 0.
        \]
        Разделив обе части на $2$, получаем:
        \[
            \sum_{\substack{k=0 \\ k \text{ чётные}}}^{n} C_n^k = 
            \sum_{k=0}^{\lfloor n/2 \rfloor} C_n^{2k} = \frac{2^n}{2} = 2^{n-1}.
        \]
        Вычитая из всей суммы $2^n$ сумму по чётным индексам $2^{n-1}$, 
        находим сумму по нечётным индексам: $2^n - 2^{n-1} = 2^{n-1}$.
    \end{quote}
\end{proof}

\subsubsection{Свойство 5. Дифференцирование бинома (взвешенная сумма)}

\begin{proposition}
    Для любого натурального $n \ge 1$ выполняется тождество:
    \begin{equation} \label{eq:binom_sum_k}
        \sum_{k=0}^{n} k C_n^k = n \cdot 2^{n-1} \tag{7}
    \end{equation}
\end{proposition}

\begin{proof}[Доказательство]~\\
    \begin{quote}
        Запишем формулу бинома Ньютона \eqref{eq:newton_binom} для 
        функции одной переменной $(1 + x)^n$:
        \[
            (1 + x)^n = \sum_{k=0}^{n} C_n^k x^k.
        \]
        Продифференцируем обе части этого тождества по переменной $x$, 
        используя формулу производной степени $(x^k)' = k x^{k-1}$:
        \[
            \bigl((1 + x)^n\bigr)' = n(1 + x)^{n-1} = 
            \sum_{k=0}^{n} k C_n^k x^{k-1}.
        \]
        Подставим в полученное равенство значение $x = 1$:
        \[
            n(1 + 1)^{n-1} = \sum_{k=0}^{n} k C_n^k \cdot 1^{k-1} \implies 
            \sum_{k=0}^{n} k C_n^k = n \cdot 2^{n-1}.
        \]
    \end{quote}
\end{proof}

\subsubsection{Свойство 6. Интегрирование бинома}

\begin{proposition}
    Для любого $n \in \N \cup \{0\}$ справедливо равенство:
    \begin{equation} \label{eq:binom_sum_int}
        \sum_{k=0}^{n} \frac{1}{k + 1} C_n^k = \frac{2^{n+1} - 1}{n + 1} \tag{8}
    \end{equation}
\end{proposition}

\begin{proof}[Доказательство]~\\
    \begin{quote}
        Рассмотрим биномиальное разложение $(1 + x)^n = \sum_{k=0}^n C_n^k x^k$ 
        и проинтегрируем обе его части по отрезку $[0, 1]$:
        \[
            \int_{0}^{1} (1 + x)^n \, dx = 
            \int_{0}^{1} \left( \sum_{k=0}^{n} C_n^k x^k \right) dx.
        \]
        В силу линейности определённого интеграла проинтегрируем сумму 
        почленно и применим формулу Ньютона---Лейбница:
        \[
            \left. \frac{(1 + x)^{n+1}}{n + 1} \right|_{0}^{1} = 
            \sum_{k=0}^{n} C_n^k \int_{0}^{1} x^k \, dx = 
            \sum_{k=0}^{n} C_n^k \left. \frac{x^{k+1}}{k + 1} \right|_{0}^{1}.
        \]
        Подставляя верхний предел $x = 1$ и нижний предел $x = 0$, получаем:
        \[
            \frac{2^{n+1} - 1^{n+1}}{n + 1} = 
            \sum_{k=0}^{n} C_n^k \cdot \frac{1}{k + 1},
        \]
        что в точности совпадает с формулой \eqref{eq:binom_sum_int}.
    \end{quote}
\end{proof}

\subsubsection{Свойство 7. Правило (тождество) Паскаля}

\begin{theorem}[\textbf{Правило Паскаля}]
    Для любых натуральных $n \ge 1$ и $1 \le k \le n$ справедливо 
    рекуррентное соотношение:
    \begin{equation} \label{eq:pascal_id}
        C_n^k = C_{n-1}^{k-1} + C_{n-1}^k \tag{9}
    \end{equation}
\end{theorem}

\begin{proof}[Доказательство (алгебраическое и комбинаторное)]~\\
    \begin{quote}
        \textbf{1. Алгебраическое доказательство:} 
        Запишем правую часть по определению через факториалы и приведём 
        дроби к общему знаменателю $k!(n - k)!$:
        \begin{align*}
            C_{n-1}^{k-1} + C_{n-1}^k 
            &= \frac{(n - 1)!}{(k - 1)!(n - k)!} + 
            \frac{(n - 1)!}{k!(n - k - 1)!} \\
            &= \frac{k \cdot (n - 1)! + (n - k) \cdot (n - 1)!}{k!(n - k)!} \\
            &= \frac{(k + n - k) \cdot (n - 1)!}{k!(n - k)!} = 
            \frac{n \cdot (n - 1)!}{k!(n - k)!} = \frac{n!}{k!(n - k)!} = C_n^k.
        \end{align*}

        \textbf{2. Комбинаторное доказательство:} 
        Выделим в $n$-элементном множестве $X = \{x_1, \dots, x_n\}$ 
        один фиксированный элемент $x_n$. Все $k$-элементные подмножества 
        множества $X$ (общее число которых равно $C_n^k$) распадаются на 
        два непересекающихся класса:
        \begin{itemize}
            \item подмножества, \emph{содержащие} элемент $x_n$: для их 
            формирования нужно добрать ещё $k - 1$ элемент из оставшихся 
            $n - 1$ элементов, что даёт $C_{n-1}^{k-1}$ вариантов;
            \item подмножества, \emph{не содержащие} элемент $x_n$: все $k$ 
            элементов выбираются из оставшихся $n - 1$ элементов, что даёт 
            $C_{n-1}^k$ вариантов.
        \end{itemize}
        По принципу сложения $C_n^k = C_{n-1}^{k-1} + C_{n-1}^k$.
    \end{quote}
\end{proof}

\subsubsection{Свойство 8. Свёртка Вандермонда}

\begin{theorem}[\textbf{Свёртка Вандермонда}]
    Для любых целых неотрицательных $n, m$ и $k \le \min(n, m)$ (или 
    произвольного $k$ при соглашении $C_r^s = 0$ для $s > r$) выполняется 
    тождество:
    \begin{equation} \label{eq:vandermonde}
        \underbrace{\sum_{p=0}^{k} C_n^p \cdot C_m^{k-p}}_{\text{свёртка}} = 
        C_{n+m}^k \tag{10}
    \end{equation}
\end{theorem}

\begin{proof}[Доказательство]~\\
    \begin{quote}
        \textbf{1. Алгебраический вывод:} 
        Рассмотрим очевидное степенное тождество:
        \[
            (1 + x)^n \cdot (1 + x)^m = (1 + x)^{n+m}.
        \]
        Разложим каждый из множителей по формуле бинома Ньютона 
        \eqref{eq:newton_binom}:
        \[
            \left(\sum_{p=0}^{n} C_n^p x^p\right) \cdot 
            \left(\sum_{q=0}^{m} C_m^q x^q\right) = 
            \sum_{k=0}^{n+m} C_{n+m}^k x^k.
        \]
        Найдём коэффициент при степени $x^k$ в обеих частях равенства. В 
        правой части он равен $C_{n+m}^k$. В левой части при перемножении 
        двух сумм степень $x^k$ возникает тогда и только тогда, когда 
        перемножаются слагаемые $C_n^p x^p$ и $C_m^q x^q$ с условием 
        $p + q = k$, то есть $q = k - p$. Суммируя по всем возможным 
        $p = 0, 1, \dots, k$, получаем коэффициент при $x^k$ в левой части:
        \[
            \sum_{p=0}^{k} C_n^p \cdot C_m^{k-p}.
        \]
        Так как два многочлена тождественно равны тогда и только тогда, 
        когда равны их коэффициенты при одинаковых степенях $x$, 
        получаем \eqref{eq:vandermonde}.

        \textbf{2. Комбинаторный вывод:} 
        Пусть имеется группа из $n + m$ объектов (например, $n$ белых и 
        $m$ чёрных шаров). Выбрать $k$ шаров из общего числа можно 
        $C_{n+m}^k$ способами. С другой стороны, в выбранном наборе может 
        оказаться ровно $p$ белых шаров и $k - p$ чёрных шаров 
        ($0 \le p \le k$). По принципу умножения и принципу сложения число 
        способов равно $\sum_{p=0}^k C_n^p C_m^{k-p}$.
    \end{quote}
\end{proof}

\subsubsection{Свойство 9. Суммирование по верхнему индексу}

\begin{proposition}[\textbf{Тождество «хоккейной клюшки»}]
    Для любых натуральных $n, m \in \N$ справедливо равенство:
    \begin{equation} \label{eq:hockey_stick}
        \sum_{k=0}^{m-1} C_{n+k}^n = C_{n+m}^{n+1} \tag{11}
    \end{equation}
\end{proposition}

\begin{proof}[Доказательство методом математической индукции]~\\
    \begin{quote}
        Проведём индукцию по числу слагаемых $m$ при фиксированном $n$.

        \textbf{1. База индукции ($m = 1$):} 
        При $m = 1$ левая часть содержит одно слагаемое при $k = 0$:
        \[
            \sum_{k=0}^{0} C_{n+k}^n = C_n^n = 1 = C_{n+1}^{n+1}.
        \]
        База индукции верна.

        \textbf{2. Индукционный шаг ($m - 1 \implies m$):} 
        Пусть тождество \eqref{eq:hockey_stick} верно для $m - 1$, то есть 
        $\sum_{k=0}^{m-2} C_{n+k}^n = C_{n+m-1}^{n+1}$. Докажем его для $m$. 
        Отделим последнее слагаемое суммы при $k = m - 1$ и воспользуемся 
        предположением индукции, а затем правилом Паскаля 
        \eqref{eq:pascal_id}:
        \begin{align*}
            \sum_{k=0}^{m-1} C_{n+k}^n 
            &= \left(\sum_{k=0}^{m-2} C_{n+k}^n\right) + C_{n+m-1}^n \\
            &= C_{n+m-1}^{n+1} + C_{n+m-1}^n = C_{n+m}^{n+1}.
        \end{align*}
        Утверждение полностью доказано.
    \end{quote}
\end{proof}

\subsection{Треугольник Паскаля и его скрытые закономерности}

\begin{definition}
    \textdef{Треугольником Паскаля} называется бесконечная треугольная 
    таблица биномиальных коэффициентов, в которой в $n$-й строке 
    ($n = 0, 1, 2, \dots$) стоят числа $C_n^0, C_n^1, \dots, C_n^n$. 
    Крайние элементы каждой строки равны $1$, а каждое внутреннее число 
    равно сумме двух расположенных над ним чисел предыдущей строки 
    согласно правилу Паскаля \eqref{eq:pascal_id}.
\end{definition}

\begin{figure}[H]
    \centering
    \begin{tikzpicture}[scale=0.95]
        % Вершина C_0^0
        \node (r0k0) at (0, 0) {$1$};
        \node[above right=2pt, gray, scale=0.85] at (0.1, 0.1) {$C_0^0$};
        \draw[->, gray, thin] (0.3, 0.25) -- (0.12, 0.12);

        % Строка 1
        \node (r1k0) at (-0.7, -0.9) {$1$};
        \node (r1k1) at (0.7, -0.9) {$1$};

        % Строка 2
        \node (r2k0) at (-1.4, -1.8) {$1$};
        \node (r2k1) at (0, -1.8) {$2$};
        \node (r2k2) at (1.4, -1.8) {$1$};

        % Строка 3
        \node (r3k0) at (-2.1, -2.7) {$1$};
        \node (r3k1) at (-0.7, -2.7) {$3$};
        \node (r3k2) at (0.7, -2.7) {$3$};
        \node (r3k3) at (2.1, -2.7) {$1$};

        % Строка 4
        \node (r4k0) at (-2.8, -3.6) {$1$};
        \node (r4k1) at (-1.4, -3.6) {$4$};
        \node (r4k2) at (0, -3.6) {$6$};
        \node (r4k3) at (1.4, -3.6) {$4$};
        \node (r4k4) at (2.8, -3.6) {$1$};

        % Пунктирные связи сложения по правилу Паскаля
        \draw[dashed, gray] (r0k0) -- (r1k0);
        \draw[dashed, gray] (r0k0) -- (r1k1);

        \draw[dashed, gray] (r1k0) -- (r2k0);
        \draw[dashed, gray] (r1k0) -- (r2k1);
        \draw[dashed, gray] (r1k1) -- (r2k1);
        \draw[dashed, gray] (r1k1) -- (r2k2);

        \draw[dashed, gray] (r2k0) -- (r3k0);
        \draw[dashed, gray] (r2k0) -- (r3k1);
        \draw[dashed, gray] (r2k1) -- (r3k1);
        \draw[dashed, gray] (r2k1) -- (r3k2);
        \draw[dashed, gray] (r2k2) -- (r3k2);
        \draw[dashed, gray] (r2k2) -- (r3k3);

        \draw[dashed, gray] (r3k0) -- (r4k0);
        \draw[dashed, gray] (r3k0) -- (r4k1);
        \draw[dashed, gray] (r3k1) -- (r4k1);
        \draw[dashed, gray] (r3k1) -- (r4k2);
        \draw[dashed, gray] (r3k2) -- (r4k2);
        \draw[dashed, gray] (r3k2) -- (r4k3);
        \draw[dashed, gray] (r3k3) -- (r4k3);
        \draw[dashed, gray] (r3k3) -- (r4k4);
    \end{tikzpicture}
    \caption{Классическое представление треугольника Паскаля и схема 
    формирования элементов по правилу Паскаля.}
\end{figure}

\subsubsection{Диагонали треугольника Паскаля и числа Фибоначчи}

Выровняем строки треугольника Паскаля по левому краю (так, чтобы элементы 
$C_n^k$ стояли в $n$-й строке и $k$-м столбце прямоугольной таблицы) и 
просуммируем числа вдоль восходящих диагоналей.

\begin{figure}[H]
    \centering
    \begin{tikzpicture}[scale=0.85]
        % Верхняя строка с суммами (числами Фибоначчи)
        \draw[thick] (-0.5, 0.7) -- (6.5, 0.7);
        \node[blue!80!black, font=\bfseries] at (0.3, 1.05) {$1$};
        \node[blue!80!black, font=\bfseries] at (1.0, 1.05) {$1$};
        \node[blue!80!black, font=\bfseries] at (1.7, 1.05) {$2$};
        \node[blue!80!black, font=\bfseries] at (2.4, 1.05) {$3$};
        \node[blue!80!black, font=\bfseries] at (3.1, 1.05) {$5$};
        \node[blue!80!black, font=\bfseries] at (3.8, 1.05) {$8$};
        \node[blue!80!black, font=\bfseries] at (4.6, 1.05) {$13$};
        \node[right, blue!80!black] at (5.3, 1.05) 
            {$\longleftarrow$ числа Фибоначчи};

        % Выровненный по левому краю треугольник Паскаля
        % Строка 0
        \node at (0, 0) {$1$};
        % Строка 1
        \node at (0, -0.7) {$1$}; \node at (0.7, -0.7) {$1$};
        % Строка 2
        \node at (0, -1.4) {$1$}; \node at (0.7, -1.4) {$2$}; 
        \node at (1.4, -1.4) {$1$};
        % Строка 3
        \node at (0, -2.1) {$1$}; \node at (0.7, -2.1) {$3$}; 
        \node at (1.4, -2.1) {$3$}; \node at (2.1, -2.1) {$1$};
        % Строка 4
        \node at (0, -2.8) {$1$}; \node at (0.7, -2.8) {$4$}; 
        \node at (1.4, -2.8) {$6$}; \node at (2.1, -2.8) {$4$}; 
        \node at (2.8, -2.8) {$1$};
        % Строка 5
        \node at (0, -3.5) {$1$}; \node at (0.7, -3.5) {$5$}; 
        \node at (1.4, -3.5) {$10$}; \node at (2.1, -3.5) {$10$}; 
        \node at (2.8, -3.5) {$5$}; \node at (3.5, -3.5) {$1$};

        % Восходящие диагонали
        \draw[blue!60, thin] (-0.2, -0.2) -- (0.3, 0.7);
        \draw[blue!60, thin] (-0.2, -0.9) -- (1.0, 0.7);
        \draw[blue!60, thin] (-0.2, -1.6) -- (1.7, 0.7);
        \draw[blue!60, thin] (-0.2, -2.3) -- (2.4, 0.7);
        \draw[blue!60, thin] (-0.2, -3.0) -- (3.1, 0.7);
        \draw[blue!60, thin] (-0.2, -3.7) -- (3.8, 0.7);
        \draw[blue!60, thin] (0.5, -3.7) -- (4.6, 0.7);
    \end{tikzpicture}
    \caption{Суммирование элементов выровненного треугольника Паскаля вдоль 
    восходящих диагоналей порождает числа Фибоначчи.}
\end{figure}

\begin{theorem}[\textbf{Связь треугольника Паскаля с числами Фибоначчи}]
    Пусть $\{F_m\}_{m=1}^\infty$ --- последовательность чисел Фибоначчи, 
    заданная условиями $F_1 = 1$, $F_2 = 1$ и $F_{m+1} = F_m + F_{m-1}$ 
    при $m \ge 2$. Тогда для любого $n \in \N \cup \{0\}$ сумма элементов 
    $n$-й восходящей диагонали треугольника Паскаля равна $(n + 1)$-му 
    числу Фибоначчи:
    \begin{equation} \label{eq:fib_binom}
        F_{n+1} = \sum_{k=0}^{\lfloor n/2 \rfloor} C_{n-k}^k \tag{12}
    \end{equation}
\end{theorem}

\begin{proof}[Доказательство]~\\
    \begin{quote}
        Обозначим сумму $n$-й восходящей диагонали через 
        $S_n = \sum_{k \ge 0} C_{n-k}^k$ (считая $C_r^s = 0$ при $s > r$ 
        или $s < 0$). Докажем, что последовательность $S_n$ удовлетворяет 
        тем же начальным условиям и тому же рекуррентному соотношению, что 
        и числа Фибоначчи $F_{n+1}$.

        \textbf{1. Проверка начальных значений ($n = 0$ и $n = 1$):}
        \[
            S_0 = C_0^0 = 1 = F_1, \qquad
            S_1 = C_1^0 = 1 = F_2.
        \]

        \textbf{2. Вывод рекуррентного соотношения ($S_{n} = S_{n-1} + S_{n-2}$):} 
        Пусть $n \ge 2$. Применим к каждому слагаемому $C_{n-k}^k$ (при 
        $k \ge 1$) правило Паскаля \eqref{eq:pascal_id}:
        \[
            C_{n-k}^k = C_{n-1-k}^k + C_{n-1-k}^{k-1}.
        \]
        Тогда вся диагональная сумма $S_n$ распадается на две суммы:
        \begin{align*}
            S_n &= C_n^0 + \sum_{k \ge 1} C_{n-k}^k = 
            C_{n-1}^0 + \sum_{k \ge 1} \bigl(C_{n-1-k}^k + C_{n-1-k}^{k-1}\bigr) \\
            &= \underbrace{\sum_{k \ge 0} C_{(n-1)-k}^k}_{S_{n-1}} + 
            \sum_{k \ge 1} C_{(n-2)-(k-1)}^{k-1}.
        \end{align*}
        Производя во второй сумме замену индекса $j = k - 1 \ge 0$, 
        получаем в точности $\sum_{j \ge 0} C_{(n-2)-j}^j = S_{n-2}$. 
        Таким образом, $S_n = S_{n-1} + S_{n-2}$, откуда по индукции 
        $S_n = F_{n+1}$ для всех $n \ge 0$.
    \end{quote}
\end{proof}

\subsubsection{Арифметика треугольника Паскаля по модулю номера строки}

Рассмотрим выровненный по левому краю треугольник Паскаля и в каждой строке 
с номером $n \ge 1$ заменим внутренние биномиальные коэффициенты $C_n^k$ 
($1 \le k \le n - 1$) на их остатки от деления на номер строки $n$ 
(крайние единицы $C_n^0 = C_n^n = 1$ при $n \ge 2$ имеют остаток $1$, а при 
$n = 1$ дают $0$).

\[
\begin{array}{r|cccccccccccccccccc}
    n & \multicolumn{8}{l}{\text{Исходный треугольник } C_n^k} & \quad & 
    \multicolumn{8}{l}{\text{Остатки по модулю } n \ (C_n^k \bmod n)} \\[1mm]
    \hline
    0: & 1 & & & & & & & & & 1 & & & & & & & \\
    1: & 1 & 1 & & & & & & & & 0 & 0 & & & & & & \\
    2: & 1 & 2 & 1 & & & & & & & 1 & \mathbf{0} & 1 & & & & & \\
    3: & 1 & 3 & 3 & 1 & & & & & & 1 & \mathbf{0} & \mathbf{0} & 1 & & & & \\
    4: & 1 & 4 & 6 & 4 & 1 & & & & & 1 & 0 & 2 & 0 & 1 & & & \\
    5: & 1 & 5 & 10 & 10 & 5 & 1 & & & & 
    1 & \mathbf{0} & \mathbf{0} & \mathbf{0} & \mathbf{0} & 1 & & \\
    6: & 1 & 6 & 15 & 20 & 15 & 6 & 1 & & & 
    1 & 0 & 3 & 2 & 3 & 0 & 1 & \\
    7: & 1 & 7 & 21 & 35 & 35 & 21 & 7 & 1 & & 
    1 & \mathbf{0} & \mathbf{0} & \mathbf{0} & \mathbf{0} & \mathbf{0} & \mathbf{0} & 1
\end{array}
\]

Если теперь каждую $n$-ю строку таблицы остатков сдвинуть вправо на $2n$ 
позиций (так, чтобы строка $n$ начиналась под столбцом с номером $2n$), 
возникает наглядная диаграмма распределения простых и составных чисел:

\begin{figure}[H]
    \centering
    \begin{tikzpicture}[scale=0.56, every node/.style={scale=0.8}]
        % Шкала номеров столбцов сверху
        \foreach \x in {0,...,21} {
            \node[gray] at (\x, 1.2) {$\x$};
        }

        % Вертикальные красные линии под простыми числами (2, 3, 5, 7, 11, 13, 17, 19)
        \foreach \p/\ymin in {2/-1, 3/-1, 5/-2, 7/-3, 11/-5, 13/-6, 17/-7, 19/-7} {
            \draw[red!70!black, thin] (\p, 0.7) -- (\p, \ymin + 0.3);
        }
        \node[right, red!70!black] at (21.6, 1.2) {простые $\to 0$};

        % Строка n = 0 (сдвиг 0)
        \node at (0, 0) {$1$};

        % Строка n = 1 (сдвиг 2)
        \node at (2, -1) {$0$}; \node at (3, -1) {$0$};

        % Строка n = 2 (сдвиг 4)
        \node at (4, -2) {$1$}; \node at (5, -2) {$0$}; 
        \node at (6, -2) {$1$};

        % Строка n = 3 (сдвиг 6)
        \node at (6, -3) {$1$}; \node at (7, -3) {$0$}; 
        \node at (8, -3) {$0$}; \node at (9, -3) {$1$};

        % Строка n = 4 (сдвиг 8)
        \node at (8, -4) {$1$}; \node at (9, -4) {$0$}; 
        \node at (10, -4) {$2$}; \node at (11, -4) {$0$}; 
        \node at (12, -4) {$1$};

        % Строка n = 5 (сдвиг 10)
        \node at (10, -5) {$1$}; \node at (11, -5) {$0$}; 
        \node at (12, -5) {$0$}; \node at (13, -5) {$0$}; 
        \node at (14, -5) {$0$}; \node at (15, -5) {$1$};

        % Строка n = 6 (сдвиг 12)
        \node at (12, -6) {$1$}; \node at (13, -6) {$0$}; 
        \node at (14, -6) {$3$}; \node at (15, -6) {$2$}; 
        \node at (16, -6) {$3$}; \node at (17, -6) {$0$}; 
        \node at (18, -6) {$1$};

        % Строка n = 7 (сдвиг 14)
        \node at (14, -7) {$1$}; \node at (15, -7) {$0$}; 
        \node at (16, -7) {$0$}; \node at (17, -7) {$0$}; 
        \node at (18, -7) {$0$}; \node at (19, -7) {$0$}; 
        \node at (20, -7) {$0$}; \node at (21, -7) {$1$};

        % Пояснение сдвига слева
        \draw[->, blue!70!black, thick] (-2.5, -1.5) -- (-0.5, -1.5);
        \node[above, blue!70!black, align=center] at (-1.5, -1.5) 
            {сдвиг на\\$2n$};
    \end{tikzpicture}
    \caption{Таблица остатков $C_n^k \bmod n$ со сдвигом каждой $n$-й 
    строки на $2n$ позиций вправо: в столбцах простых чисел стоят только 
    нули.}
\end{figure}

Обоснуем математически оба фундаментальных свойства, наблюдаемых в этой 
таблице.

\begin{theorem}[\textbf{Биномиальный критерий простоты числа}]
    Натуральное число $n \ge 2$ является простым тогда и только тогда, 
    когда все внутренние биномиальные коэффициенты $n$-й строки делятся 
    нацело на $n$:
    \begin{equation} \label{eq:prime_binom_crit}
        (n \text{ --- простое}) \iff 
        \forall k \in \{1, 2, \dots, n - 1\} : C_n^k \equiv 0 \pmod{n}. \tag{13}
    \end{equation}
\end{theorem}

\begin{proof}[Доказательство критерия простоты и свойства сдвинутых столбцов]~\\
    \begin{quote}
        \textbf{1. Необходимость ($n$ --- простое $\implies C_n^k \vdots n$):}\\
        Пусть $n = p$ --- простое число, и пусть $1 \le k \le p - 1$. 
        Запишем тождество:
        \[
            C_p^k = \frac{p!}{k!(p - k)!} \implies 
            k!(p - k)! \cdot C_p^k = p \cdot (p - 1)!.
        \]
        Правая часть равенства делится на простое число $p$, значит, и 
        левая часть $k!(p - k)! \cdot C_p^k$ делится на $p$. Но все 
        множители в факториалах $k! = 1 \cdot 2 \cdot \dots \cdot k$ и 
        $(p - k)! = 1 \cdot 2 \cdot \dots \cdot (p - k)$ строго меньше 
        простого числа $p$ (так как $1 \le k \le p - 1$). Следовательно, 
        ни один из них не делится на $p$, откуда 
        $\operatorname{НОД}\bigl(p, k!(p - k)!\bigr) = 1$. По лемме Евклида 
        на $p$ обязан делиться второй сомножитель, то есть $C_p^k \vdots p$.

        \textbf{2. Достаточность ($n$ --- составное $\implies \exists k : 
        C_n^k \not\vdots n$):}\\
        Пусть $n$ --- составное число, и пусть $p$ --- его наименьший 
        простой делитель ($1 < p < n$). Рассмотрим биномиальный коэффициент 
        $C_n^p$ и преобразуем его:
        \[
            C_n^p = \frac{n}{p} \cdot C_{n-1}^{p-1} \implies 
            \frac{C_n^p}{n} = \frac{C_{n-1}^{p-1}}{p} = 
            \frac{(n - 1)(n - 2) \dots (n - p + 1)}{p!}.
        \]
        Чтобы число $C_n^p$ делилось нацело на $n$, необходимо и 
        достаточно, чтобы число $C_{n-1}^{p-1}$ делилось нацело на $p$. 
        Однако $n$ делится на $p$ ($n \equiv 0 \pmod{p}$), поэтому в 
        числителе произведения $(n - 1)(n - 2) \dots (n - p + 1)$ ни один 
        из $p - 1$ последовательных сомножителей не делится на $p$ (они 
        сравнимы с $-1, -2, \dots, -(p - 1)$ по модулю $p$). Значит, 
        числитель не делится на простое число $p$, тогда как знаменатель 
        $p! = p \cdot (p - 1)!$ содержит простой множитель $p$. 
        Следовательно, $C_{n-1}^{p-1}$ не делится на $p$, а $C_n^p$ не 
        делится на $n$, то есть $C_n^p \not\equiv 0 \pmod{n}$.

        \textbf{3. Следствие для сдвинутой таблицы (почему в простых 
        столбцах стоят нули):}\\
        При сдвиге $n$-й строки на $2n$ позиций вправо крайние единицы 
        $C_n^0 = 1$ и $C_n^n = 1$ при $n \ge 2$ попадают в столбцы с 
        номерами $2n$ и $3n$, которые всегда являются \emph{составными} 
        числами (кратными $2$ и $3$ соответственно). Поэтому в столбец с 
        простым номером $m \ge 2$ могут попасть только \emph{внутренние} 
        элементы строк вида $C_n^{m-2n} \bmod n$ (где $1 \le m - 2n \le n - 1$, 
        то есть $\frac{m}{3} < n < \frac{m}{2}$). Для всех первых простых 
        чисел $m \in \{2, 3, 5, 7, 11, 13, 17, 19\}$ попадающие в их 
        столбцы элементы имеют остаток $0$.
    \end{quote}
\end{proof}